% Controlled expansion: expanding exactly one step, and stopping.
%
% expl3 is built out of these. `\exp:w` is \romannumeral and `\exp_end:` a
% \chardef'd zero, so `\exp:w <tokens> ` expands <tokens> while scanning a
% number and stops the moment one of them leaves \exp_end: in front. Whatever
% the expansion built is then left *unexpanded* -- which is the whole point,
% and what makes an over-eager scanner so damaging: it runs the very token the
% idiom exists to protect.
\input plain
\catcode`\:=11 \catcode`\_=11
\chardef\c_zero_int=0
\let\exp:w\romannumeral \let\exp_end:\c_zero_int
\let\exp_after:wN\expandafter
\def\exp_not:v#1{RAN(#1)}
\def\exp_args:NNc#1#2#3{\exp_after:wN#1\exp_after:wN#2\csname #3\endcsname}
% Under \edef everything expands afterwards, so both forms give the same text.
\edef\a{\exp:w \exp_args:NNc \exp_after:wN \exp_end: {exp_not:v}{ARG}}
\message{<edef=\meaning\a>}
% Under a single expansion the difference shows: the built control sequence
% must survive. `\unexpanded\exp_after:wN{...}` is expl3's way of expanding
% exactly one token inside an otherwise protected group -- so \unexpanded has
% to expand to find its own opening brace (tex.web 403).
\edef\b{\unexpanded\exp_after:wN{\exp:w \exp_args:NNc \exp_after:wN \exp_end: {exp_not:v}{ARG}}}
\message{<once=\meaning\b>}
\edef\c{\unexpanded\exp_after:wN{\exp:w \exp_after:wN \exp_end: \csname exp_not:v\endcsname{ARG}}}
\message{<csname=\meaning\c>}
% \expanded and \detokenize find their brace the same way.
\def\inner{IN}
\edef\d{\expanded\exp_after:wN{\inner}}\message{<expanded=\meaning\d>}
\edef\e{\detokenize\exp_after:wN{\inner}}\message{<detok=\meaning\e>}
\bye
